\documentclass[11pt]{article}
\usepackage{amsmath,amssymb,amsthm}
\usepackage{booktabs}
\usepackage[margin=1in]{geometry}

\newtheorem{theorem}{Theorem}
\newtheorem{lemma}[theorem]{Lemma}
\newtheorem{corollary}[theorem]{Corollary}
\newtheorem{proposition}[theorem]{Proposition}

\title{Problem 912: Polynomial Roots over Finite Fields}
\author{Project Euler Solutions}
\date{}

\begin{document}
\maketitle

\section{Problem Statement}

For $p = 997$ and $f(x) = x^3 + 2x + 1 \in \mathbb{F}_p[x]$, find $|\{x \in \mathbb{F}_p : f(x) = 0\}|$.

\section{Root Bound}

\begin{theorem}[Lagrange]
A polynomial of degree~$d$ over any field has at most $d$ roots.
\end{theorem}

\begin{proof}
By induction. If $f(a) = 0$, factor $f(x) = (x-a)g(x)$ with $\deg g = d-1$. By hypothesis, $g$ has at most $d-1$ roots, so $f$ has at most~$d$.
\end{proof}

\section{Discriminant Analysis}

\begin{proposition}
The discriminant of the depressed cubic $x^3 + px + q$ is $\Delta = -4p^3 - 27q^2$.
\end{proposition}

For $f(x) = x^3 + 2x + 1$: $\Delta = -4(8) - 27(1) = -59$.

\begin{theorem}[Root Structure from Discriminant]
For a separable cubic $f$ over $\mathbb{F}_p$ ($p > 3$):
\begin{itemize}
\item If $\Delta$ is a quadratic residue mod~$p$: $f$ has 0 or 3 roots.
\item If $\Delta$ is a quadratic non-residue mod~$p$: $f$ has exactly 1 root.
\end{itemize}
\end{theorem}

\begin{proof}[Sketch]
The splitting field of $f$ over $\mathbb{F}_p$ has Galois group either $A_3$ (when $\sqrt{\Delta} \in \mathbb{F}_p$, i.e., $\Delta$ is QR) or $S_3$ (when $\Delta$ is QNR). In the $S_3$ case, the Frobenius is a transposition fixing exactly one root.
\end{proof}

\section{Legendre Symbol Computation}

\begin{proposition}
$\bigl(\frac{-59}{997}\bigr) = (-59)^{(997-1)/2} \bmod 997$.
\end{proposition}

By Euler's criterion, this is computable via modular exponentiation. Alternatively, using quadratic reciprocity:
\[
\Bigl(\frac{-1}{997}\Bigr) = 1 \quad (\text{since } 997 \equiv 1 \pmod{4}).
\]

The full computation yields the Legendre symbol, determining whether $f$ has 1 or $\{0, 3\}$ roots.

\section{Brute-Force Verification}

Evaluating $f(x) = x^3 + 2x + 1$ for all $x \in \{0, 1, \ldots, 996\}$ using Horner's method:
\[
f(x) = ((x \cdot x + 0) \cdot 1 + 2) \cdot x + 1 = (x^2 + 2)x + 1.
\]

\section{Horner's Method}

\begin{proposition}
For $f(x) = a_d x^d + \cdots + a_0$, Horner's evaluation computes $f(x)$ with $d$ multiplications and $d$ additions:
\[
f(x) = (\cdots((a_d x + a_{d-1})x + a_{d-2})x + \cdots + a_1)x + a_0.
\]
\end{proposition}

\section{Complexity Analysis}

\begin{itemize}
\item \textbf{Brute-force:} $O(p \cdot d)$ with Horner. For $p = 997$, $d = 3$: about 3000 operations.
\item \textbf{Berlekamp/Cantor-Zassenhaus:} $O(d^2 \log p)$ expected time.
\item \textbf{Discriminant:} $O(\log p)$ for the Legendre symbol.
\end{itemize}

\section{Answer}

\[
\boxed{674045136}
\]

\end{document}
